\pagestyle{empty}
\tight
\begin{tblr}{width=\textwidth,colspec={|Q[c,m,1.9cm]|X[l,m]|},hlines,vlines,hspan=minimal,row{1}={bg=headblue},rowsep=1pt,colsep=3pt}
\SetCell{c,m}\hfil\logo\hfil & \textbf{POLITEKNIK NEGERI MALANG}\par\textbf{JURUSAN TEKNOLOGI INFORMASI}\par\textbf{PROGRAM STUDI: D4 TEKNIK INFORMATIKA} \\
\SetCell[c=2]{c} \textbf{RENCANA PEMBELAJARAN SEMESTER (RPS)} & \\
\end{tblr}
\begin{longtblr}{width=\textwidth,colspec={|Q[l,m,1.9cm]|X[0.85,l,m]|X[1.45,l,m]|X[0.75,l,m]|X[1.15,l,m]|},hlines,vlines,hspan=minimal,rowsep=1pt,colsep=2pt,row{1,3}={bg=headgray,font=\bfseries}}
MATA KULIAH & KODE & BOBOT (SKS)/jam & SEMESTER & TGL. PENYUSUNAN \\
Bahasa Inggris 1 & RTI251005 & 2 SKS/4 jam & 1 & 7 Agustus 2025 \\
OTORISASI & Dosen Pengembang RPS & Koordinator RMK & \SetCell[c=2]{l} Ka PRODI &  \\
 & Atiqah Nurul Asri, S.Pd., M.Pd. & Atiqah Nurul Asri, S.Pd., M.Pd. & \SetCell[c=2]{l}{\vspace{8mm}\par Dr. Ely Setyo Astuti, S.T., M.T.} &  \\
\SetCell[r=9]{l} Capaian Pembelajaran (CP) & \SetCell[c=4]{l,bg=headgray,font=\bfseries} Capaian Pembelajaran Lulusan Program Studi (CPL-Prodi) &  &  &  \\
 & CPL03 & \SetCell[c=3]{l} Mampu mengelola tim dan diri sendiri, berkomunikasi secara lisan dan tulisan, serta mempresentasikan gagasan dan hasil kerja secara efektif dan profesional. &  &  \\
 & \SetCell[c=4]{l,bg=headgray,font=\bfseries} Capaian Pembelajaran mata kuliah (CPL-MK) &  &  &  \\
 & CPMK0303 & \SetCell[c=3]{l} Mahasiswa mampu mengembangkan diri secara berkelanjutan, berkomunikasi secara profesional, mempresentasikan diri, dan mengelola karir. &  &  \\
 & \SetCell[c=4]{l,bg=headgray,font=\bfseries} Kemampuan akhir tiap tahapan belajar (Sub-CPMK) &  &  &  \\
 & SCPMK0303-00501 & \SetCell[c=3]{l} Mahasiswa mampu memahami dan menggunakan vocabulary serta grammatical functions bahasa Inggris (Simple Present Tense, imperatives, sequencers, adjectives, comparatives, superlatives, passive sentences, time clauses, if-clause, dan modals) dalam konteks Teknik Informatika. &  &  \\
 & SCPMK0303-00502 & \SetCell[c=3]{l} Mahasiswa mampu memahami teks lisan dan tulisan berbahasa Inggris tentang topik Teknik Informatika serta menjelaskan kembali isinya secara tepat. &  &  \\
 & SCPMK0303-00503 & \SetCell[c=3]{l} Mahasiswa mampu menulis teks fungsional bahasa Inggris (instruksi, perbandingan, deskripsi proses, review, dan deskripsi pekerjaan) dalam konteks Teknik Informatika. &  &  \\
 & SCPMK0303-00504 & \SetCell[c=3]{l} Mahasiswa mampu mempresentasikan gagasan dan hasil kerja secara lisan dalam bahasa Inggris secara mandiri, efektif, dan profesional. &  &  \\
\SetCell[r=2]{l} Penilaian & \SetCell[c=4]{l} *Nilai CPMK didapat dari agregasi nilai SUB-CPMK yang dibebankan &  &  &  \\
 & \SetCell[c=4]{l}{\tiny\begin{tblr}{width=\linewidth,colspec={|Q[c,m,0.1\linewidth]|Q[c,m,0.16\linewidth]|X[l,m]|X[l,m]|X[l,m]|X[l,m]|X[l,m]|Q[c,m,0.08\linewidth]|},hlines,vlines,hspan=minimal,rowsep=1pt,colsep=0.7pt,row{1,2}={bg=headgray,font=\bfseries}}
Kode CPMK & Kode SCPMK & \SetCell[c=5]{c} Bobot per Bentuk Penilaian & & & & & Total Bobot \\
 & & Tugas & Kuis & Praktik Lisan & UTS & UAS & \\
\SetCell[r=1]{c} \SetCell[r=4]{c} CPMK0303 & SCPMK0303-00501 & & 10\%: kuis 1 dan 2 vocabulary-grammar & & 5\%: vocabulary dan grammar & 5\%: vocabulary dan grammar & 20\% \\
\SetCell[r=3]{c}  & SCPMK0303-00502 & & & 5\%: pemahaman listening-reading dan tanya jawab & 5\%: pemahaman bacaan/simakan & 5\%: pemahaman bacaan/simakan & 15\% \\
 & SCPMK0303-00503 & 5\%: lima tugas menulis terstruktur & & 10\%: teks dan bahan presentasi & 5\%: menulis karangan & 5\%: menulis karangan & 25\% \\
 & SCPMK0303-00504 & & & 40\%: presentasi lisan, diskusi, dan role play & & & 40\% \\
\SetCell[c=7]{r}\textbf{Total Per Penilaian} & & & & & & & \textbf{100\%} \\
\end{tblr}} &  &  &  \\
Diskripsi Singkat Mata Kuliah & \SetCell[c=4]{l} Mata kuliah Bahasa Inggris 1 melatih kemampuan dan kecakapan mahasiswa dalam Listening, Speaking, Reading, dan Writing secara terintegrasi dalam konteks Teknik Informatika. Topik-topik materi disesuaikan dengan bidang Teknik Informatika yang dapat diterapkan dalam kehidupan sehari-hari dan dunia kerja. Aktivitas belajar meliputi ceramah, diskusi, role play, presentasi, serta project individu maupun kelompok, dan mendukung penerapan Case Method dan Project Based Learning. &  &  &  \\
Bahan Kajian / Materi Pembelajaran & \SetCell[c=4]{l}\begin{enumerate}\item Computer Applications: kinds of applications, Simple Present Tense, imperatives, sequencers, software installation process\item Computer Architecture: types of computers, hardware and its functions, computer ads, adjectives, comparatives, superlatives\item Multimedia: multimedia hardware, toolbox of a graphic program, passive sentences, time clauses\item Computer Network: network hardware, network topology, if-clause\item Websites: features, types, analytics, data visualization, characteristics of a good website\item Careers in IT: IT jobs and responsibilities, modals for skills and qualifications\item IT Support Staff: computer problems and solutions, service report, help desk ticket, email\item Workstation Health and Safety: problems, expressions for giving advice and declaring prohibition\end{enumerate} &  &  &  \\
\SetCell[r=2]{l} Pustaka & \SetCell[c=4]{l} \textbf{Utama:}\begin{enumerate}\item Asri, Atiqah Nurul. 2025. English for Informatics 1: 8th Edition. Modul belum dipublikasikan.\end{enumerate} &  &  &  \\
 & \SetCell[c=4]{l} \textbf{Pendukung:}\begin{enumerate}\item Esteras, Santiago Remacha. (2010). Infotech English for Computer Users Workbook. Cambridge: Cambridge University Press.\item Esteras, Santiago Remacha. (2011). Infotech English for Computer Users Teacher's Book. Cambridge: Cambridge University Press.\item Glendinning, Eric H \& McEwan, John. (2012). Basic English for Computing Revised and Updated. Oxford: Oxford University Press.\item Olejniczak, Maja. (2011). English for Information Technology 1 Vocational English Course Book. Essex: Pearson Education Limited.\end{enumerate} &  &  &  \\
Media Pembelajaran & \SetCell[c=2]{l} \textbf{Software:} LMS, Paddlet, Mentimeter. &  & \SetCell[c=2]{l} \textbf{Hardware:} Komputer, LCD Projector, Audio Files, dan Speakers. &  \\
Nama Dosen Pengampu & \SetCell[c=4]{l}\begin{enumerate}\item Atiqah Nurul Asri, S.Pd., M.Pd.\item Satrio Binusa Suryadi, S.S., M.Pd.\end{enumerate} &  &  &  \\
Matakuliah Syarat & \SetCell[c=4]{l} - &  &  &  \\
\end{longtblr}
\begin{longtblr}{width=\textwidth,colspec={|Q[c,m,0.06\textwidth]|X[1.35,l,m]|X[1.08,l,m]|X[1.16,l,m]|X[0.7,l,m]|X[1.24,l,m]|X[1.06,l,m]|X[0.98,l,m]|Q[c,m,0.05\textwidth]|},hlines,vlines,hspan=minimal,rowhead=2,row{1,2}={bg=headgreen,font=\bfseries},rowsep=1pt,colsep=1.5pt}
Mgg.\par Ke & Kemampuan Akhir Yang Direncanakan (Sub-CP-MK) & Bahan kajian (Materi Pembelajaran) & Bentuk dan Metode Pembelajaran & Estimasi Waktu & Pengalaman Belajar Mahasiswa & Kriteria \& Bentuk Penilaian & Indikator Penilaian & Bobot Penilaian (\%) \\
(1) & (2) & (3) & (4) & (5) & (6) & (7) & (8) & (9) \\
1 & SCPMK0303-00501, SCPMK0303-00502 & Topic 1: Computer Applications; kinds of computer applications and their uses; Simple Present Tense, imperatives, sequencers. & Bentuk: tatap muka, tugas terstruktur, tugas mandiri. Metode: ceramah dan tanya jawab, small group discussion, case method, presentasi. & TM: 2 x 50'; PT: 1 x 50'; M: 1 x 50'. & Mahasiswa mengidentifikasi berbagai aplikasi komputer beserta user-nya, membaca bacaan Nowadays Computer Uses, dan berlatih Simple Present Tense, imperatives, dan sequencers. & Kriteria: ketepatan dan penguasaan vocabulary-grammar. Bentuk: tes lisan tanya jawab, diskusi, dan presentasi. & Ketepatan mengidentifikasi aplikasi komputer; ketepatan penggunaan Simple Present Tense, imperatives, dan sequencers. & 4 \\
2 & SCPMK0303-00503, SCPMK0303-00504 & Topic 1: software installation process; Tugas 1: writing and presenting instructions on how to install software. & Bentuk: tatap muka, tugas terstruktur, tugas mandiri. Metode: problem-based learning, case method, presentasi. & TM: 2 x 50'; PT: 1 x 50'; M: 1 x 50'. & Mahasiswa menulis instruksi proses instalasi sebuah software dengan grammatical functions yang dipelajari lalu mempresentasikannya. & Kriteria: ketepatan komunikasi lisan-tulisan. Bentuk: tugas tertulis dan praktik lisan presentasi. Mengacu rubrik RTM. & Ketepatan Simple Present Tense, imperatives, dan sequencers dalam menulis dan mempresentasikan instruksi instalasi software. & 4 \\
3 & SCPMK0303-00501, SCPMK0303-00502 & Topic 2: Computer Architecture; types of computers; computer hardware and its functions. & Bentuk: tatap muka, tugas terstruktur, tugas mandiri. Metode: ceramah dan tanya jawab, small group discussion, case method. & TM: 2 x 50'; PT: 1 x 50'; M: 1 x 50'. & Mahasiswa mengidentifikasi jenis-jenis komputer dari gambar dan percakapan yang didengar serta menjelaskan komponen komputer dan fungsinya. & Kriteria: ketepatan dan penguasaan vocabulary-grammar. Bentuk: tes lisan tanya jawab, diskusi, dan presentasi. & Ketepatan membedakan jenis komputer; ketepatan menjelaskan fungsi computer hardware dengan struktur kalimat yang tepat. & 4 \\
4 & SCPMK0303-00503, SCPMK0303-00504 & Topic 2: computer ads; adjectives, comparatives, superlatives; Tugas 2: writing and presenting the comparison of devices. & Bentuk: tatap muka, tugas terstruktur, tugas mandiri. Metode: problem-based learning, case method, presentasi. & TM: 2 x 50'; PT: 1 x 50'; M: 1 x 50'. & Mahasiswa membaca iklan komputer, menuliskan perbandingan spesifikasi dua atau lebih device, dan mempresentasikannya. & Kriteria: ketepatan komunikasi lisan-tulisan. Bentuk: tugas tertulis dan praktik lisan presentasi. Mengacu rubrik RTM. & Ketepatan adjectives, comparatives, dan superlatives dalam menulis dan mempresentasikan perbandingan device. & 4 \\
5 & SCPMK0303-00501, SCPMK0303-00502 & Topic 3: Multimedia; multimedia hardware; toolbox of a graphic program. & Bentuk: tatap muka, tugas terstruktur, tugas mandiri. Metode: ceramah dan tanya jawab, small group discussion, case method. & TM: 2 x 50'; PT: 1 x 50'; M: 1 x 50'. & Mahasiswa mengidentifikasi kegunaan multimedia equipment, membaca bacaan Input Devices as Multimedia Equipment, dan menyebutkan fungsi toolbox program grafis. & Kriteria: ketepatan dan penguasaan vocabulary-grammar. Bentuk: tes lisan tanya jawab, diskusi, dan presentasi. & Ketepatan menjelaskan kegunaan multimedia equipment dan fungsi toolbox; ketepatan menjawab reading comprehension. & 4 \\
6 & SCPMK0303-00501 & Topic 3: passive sentences dan time clauses; Kuis 1 vocabulary dan grammar Topic 1--3. & Bentuk: tatap muka, tugas terstruktur, tugas mandiri. Metode: ceramah dan tanya jawab, latihan terbimbing, kuis. & TM: 2 x 50'; PT: 1 x 50'; M: 1 x 50'. & Mahasiswa berlatih passive sentences dan time clauses dalam kalimat/paragraf lalu mengerjakan Kuis 1. & Kriteria: ketepatan vocabulary dan grammar. Bentuk: kuis tertulis. Mengacu rubrik RTM. & Ketepatan penggunaan passive sentences, time clauses, dan vocabulary Topic 1--3 dalam kalimat/paragraf. & 5 \\
7 & SCPMK0303-00503, SCPMK0303-00504 & Topic 3: proses editing gambar/video/audio; Tugas 3: writing and describing an editing process. & Bentuk: tatap muka, tugas terstruktur, tugas mandiri. Metode: problem-based learning, case method, presentasi. & TM: 2 x 50'; PT: 1 x 50'; M: 1 x 50'. & Mahasiswa menulis proses editing gambar/video/audio menggunakan passive sentences dan time clauses lalu mempresentasikannya. & Kriteria: ketepatan komunikasi lisan-tulisan. Bentuk: tugas tertulis dan praktik lisan presentasi. Mengacu rubrik RTM. & Ketepatan passive sentences, time clauses, dan grammatical functions Topic 1--2 dalam menulis dan mempresentasikan proses editing. & 5 \\
8 & SCPMK0303-00501, SCPMK0303-00502, SCPMK0303-00503 & UTS: materi Topic 1--3. & Ujian tertulis dan/atau lisan. & TM: 2 x 50'. & Mahasiswa mengerjakan soal vocabulary-grammar, reading/listening comprehension, dan menulis karangan pendek. & Kriteria: ketepatan dan penguasaan komunikasi bahasa Inggris. Bentuk: UTS. Mengacu rubrik RTM. & Ketepatan vocabulary-grammar; pemahaman bacaan/simakan; kualitas karangan. & 15 \\
9 & SCPMK0303-00501, SCPMK0303-00502 & Topic 4: Computer Network; network dan network hardware; network topology; if-clause type 1. & Bentuk: tatap muka, tugas terstruktur, tugas mandiri. Metode: ceramah dan tanya jawab, small group discussion, case method. & TM: 2 x 50'; PT: 1 x 50'; M: 1 x 50'. & Mahasiswa menjelaskan pengertian network, komponen hardware network, membaca bacaan Networks dan Network Topology, serta berlatih if-clause type 1. & Kriteria: ketepatan dan penguasaan vocabulary-grammar. Bentuk: tes lisan tanya jawab, diskusi, dan presentasi. & Ketepatan menjelaskan network hardware dan topology; ketepatan penggunaan if-clause type 1 dalam kalimat/paragraf. & 4 \\
10 & SCPMK0303-00501, SCPMK0303-00502 & Topic 5: Websites; website features and their functions; types of websites and their purposes. & Bentuk: tatap muka, tugas terstruktur, tugas mandiri. Metode: ceramah dan tanya jawab, small group discussion, case method. & TM: 2 x 50'; PT: 1 x 50'; M: 1 x 50'. & Mahasiswa mengidentifikasi fitur-fitur website, membaca bacaan The World Wide Web, dan menjelaskan tipe-tipe website beserta tujuannya. & Kriteria: ketepatan dan penguasaan vocabulary-grammar. Bentuk: tes lisan tanya jawab, diskusi, dan presentasi. & Ketepatan mengidentifikasi fitur dan tipe website; ketepatan menjawab reading comprehension. & 4 \\
11 & SCPMK0303-00503, SCPMK0303-00504 & Topic 5: website analytics dan data visualization; characteristics of a good website; Tugas 4: writing and presenting a website review. & Bentuk: tatap muka, tugas terstruktur, tugas mandiri. Metode: problem-based learning, case method, presentasi. & TM: 2 x 50'; PT: 1 x 50'; M: 1 x 50'. & Mahasiswa menganalisis sebuah website berdasarkan kriteria website yang baik, menjelaskan data pada chart, menulis review, dan mempresentasikannya. & Kriteria: ketepatan komunikasi lisan-tulisan. Bentuk: tugas tertulis dan praktik lisan presentasi. Mengacu rubrik RTM. & Ketepatan adjectives, Simple Present Tense, passive voice, dan prepositions dalam menulis dan mempresentasikan review website. & 5 \\
12 & SCPMK0303-00501, SCPMK0303-00502 & Topic 6: Careers in IT; IT jobs and their responsibilities; modals for expressing skills and qualifications. & Bentuk: tatap muka, tugas terstruktur, tugas mandiri. Metode: ceramah dan tanya jawab, small group discussion, case method. & TM: 2 x 50'; PT: 1 x 50'; M: 1 x 50'. & Mahasiswa menyebutkan pekerjaan bidang IT beserta deskripsi dan kualifikasinya, membaca bacaan Industry Overview on IT Careers, dan berlatih modals. & Kriteria: ketepatan dan penguasaan vocabulary-grammar. Bentuk: tes lisan tanya jawab, diskusi, dan presentasi. & Ketepatan menjelaskan IT jobs dan kualifikasinya; ketepatan penggunaan modals dalam kalimat/paragraf. & 4 \\
13 & SCPMK0303-00503, SCPMK0303-00504 & Topic 6: deskripsi pekerjaan impian; Tugas 5: writing and talking about a dream job (video dengan subtitle). & Bentuk: tatap muka, tugas terstruktur, tugas mandiri. Metode: problem-based learning, case method, presentasi. & TM: 2 x 50'; PT: 1 x 50'; M: 1 x 50'. & Mahasiswa menulis deskripsi pekerjaan yang diimpikan menggunakan modals lalu mempresentasikannya dalam bentuk short video dengan subtitle. & Kriteria: ketepatan komunikasi lisan-tulisan. Bentuk: tugas tertulis dan praktik lisan presentasi video. Mengacu rubrik RTM. & Ketepatan modals dalam menulis dan mendeskripsikan pekerjaan impian secara lisan dan tulisan. & 5 \\
14 & SCPMK0303-00501, SCPMK0303-00502 & Topic 7: IT Support Staff; computer problems and their solutions; modals untuk analisis masalah; service report, help desk ticket, email; Kuis 2. & Bentuk: tatap muka, tugas terstruktur, tugas mandiri. Metode: ceramah dan tanya jawab, role play, case method, kuis. & TM: 2 x 50'; PT: 1 x 50'; M: 1 x 50'. & Mahasiswa mengidentifikasi masalah komputer dan solusinya, mempraktikkan percakapan via telephone, menulis email/service report, dan mengerjakan Kuis 2. & Kriteria: ketepatan dan penguasaan komunikasi. Bentuk: kuis tertulis serta tes lisan role play dan diskusi. Mengacu rubrik RTM. & Ketepatan modals dalam memberi solusi masalah komputer; kelengkapan service report/help desk ticket/email; ketepatan vocabulary-grammar Kuis 2. & 9 \\
15 & SCPMK0303-00501, SCPMK0303-00502 & Topic 8: Workstation Health and Safety; work health and safety problems; expressions for giving advice and declaring prohibition. & Bentuk: tatap muka, tugas terstruktur, tugas mandiri. Metode: ceramah dan tanya jawab, small group discussion, case method. & TM: 2 x 50'; PT: 1 x 50'; M: 1 x 50'. & Mahasiswa mengidentifikasi masalah keselamatan penggunaan lab komputer, membaca bacaan Computer Ergonomics, dan memberikan solusi atau tindakan pencegahan. & Kriteria: ketepatan dan penguasaan vocabulary-grammar. Bentuk: tes lisan tanya jawab, diskusi, dan presentasi. & Ketepatan mengidentifikasi masalah work health and safety beserta solusinya; ketepatan menjelaskan kembali isi bacaan. & 4.5 \\
16 & SCPMK0303-00503, SCPMK0303-00504 & Topic 8: peraturan penggunaan laboratorium komputer; poster slogan work health and safety. & Bentuk: tatap muka, tugas terstruktur, tugas mandiri. Metode: problem-based learning, case method, presentasi. & TM: 2 x 50'; PT: 1 x 50'; M: 1 x 50'. & Mahasiswa membuat peraturan penggunaan komputer dan laboratorium, mendesain poster slogan work health and safety dengan modals, dan mempresentasikannya. & Kriteria: ketepatan komunikasi lisan-tulisan. Bentuk: praktik lisan presentasi dan artefak poster. Mengacu rubrik RTM. & Ketepatan expressions for giving advice and declaring prohibition serta modals dalam peraturan dan poster. & 4.5 \\
17 & SCPMK0303-00501, SCPMK0303-00502, SCPMK0303-00503 & UAS: materi Topic 4--8. & Ujian tertulis dan/atau lisan. & TM: 2 x 50'. & Mahasiswa mengerjakan soal vocabulary-grammar, reading/listening comprehension, dan menulis karangan pendek. & Kriteria: ketepatan dan penguasaan komunikasi bahasa Inggris. Bentuk: UAS. Mengacu rubrik RTM. & Ketepatan vocabulary-grammar; pemahaman bacaan/simakan; kualitas karangan. & 15 \\
\end{longtblr}
